\documentclass[11pt]{article}

\usepackage[margin=1in]{geometry}
\usepackage{amsmath,amssymb,mathtools}
\usepackage{algorithm}
\usepackage{algpseudocode}
\usepackage[T1]{fontenc}
\usepackage{lmodern}
\usepackage{microtype}

\newcommand{\R}{\mathbb{R}}
\newcommand{\E}{\mathbb{E}}
\newcommand{\clip}{\mathrm{clip}}
\newcommand{\Reg}{\mathrm{Reg}}

\title{Implementation Note: Decision-Weight Regression via Cross-Fitted Alpha Weights}
\date{}

\begin{document}
\maketitle

\paragraph{Problem setting.}
For each context $z_i \in \R^p$, we observe one realized cost vector
$c_i \in \R^d$ and evaluate decisions through the linear optimization problem
\[
  x^*(c) \in \arg\min_{x \in \mathcal{X}} c^\top x.
\]
Given a predicted cost vector $\hat c_i$, the realized decision regret is
\[
  \Reg_i
  := c_i^\top\!\bigl(x^*(\hat c_i) - x^*(c_i)\bigr)
  \ge 0.
\]

\paragraph{Goal.}
We want to replace noisy label-dependent decision weights by context-only
weights $\alpha_i \approx \alpha(z_i)$, then use those $\alpha_i$ as ordinary
sample weights in a weighted least-squares cost fit.

\section*{1. Pilot model and regret-based raw weights}

We first fit an ordinary least-squares pilot model
\[
  \hat f \in \arg\min_{f \in \mathcal{F}}
  \sum_{i=1}^n \|f(z_i) - c_i\|_2^2.
\]
Using the predicted costs $\hat c_i := \hat f(z_i)$, we compute the realized
decision regret $\Reg_i$.

The implementation then forms bounded raw decision weights
\[
  w_i^{\mathrm{raw}}
  := (1-\mu) + \mu \, \phi(\Reg_i),
  \qquad \mu \in [0,1].
\]
Here $\phi$ is a bounded monotone transform. In the current implementation,
\[
  \phi(r) := \clip_{[0,1]}\!\left(\frac{r}{q_{0.9}}\right),
\]
where $q_{0.9}$ is the empirical $90$th percentile of the training regrets.
If the percentile is zero, the code falls back to the maximum regret, and if
that is also zero then all transformed regrets are set to zero.

Therefore the raw weights satisfy
\[
  w_i^{\mathrm{raw}} \in [1-\mu,\,1-\mu+\mu W_{\max}],
\]
with the current default $W_{\max}=1$.

\section*{2. Alpha-weight regression}

The key step is to estimate context-only weights by regressing these raw
decision weights on context.

Let $I_1,\dots,I_M$ be folds of $\{1,\dots,n\}$. For each fold $m$, we fit a
weight regression model on the complement:
\[
  q^{(-m)}
  \in \arg\min_{q \in \mathcal{H}}
  \sum_{i \notin I_m} \bigl(w_i^{\mathrm{raw}} - q(z_i)\bigr)^2.
\]
Then for each held-out example $i \in I_m$, we define the alpha weight
\[
  \hat\alpha_i
  := \clip_{[1-\mu,\,1-\mu+\mu W_{\max}]}\!\bigl(q^{(-m)}(z_i)\bigr).
\]

\paragraph{Interpretation.}
The regression target is the raw decision weight, which is itself a transformed
regret signal. Thus the learned alpha weights are context-smoothed versions of
the regret-based weighting rule:
\[
  z_i \mapsto \hat\alpha_i
  \approx \E\!\left[w^{\mathrm{raw}} \mid Z = z_i\right].
\]

\section*{3. Weighted least-squares refit}

Once the out-of-fold alpha weights are available, we fit the final cost model
with weighted squared loss:
\[
  \tilde f
  \in \arg\min_{f \in \mathcal{F}}
  \sum_{i=1}^n \hat\alpha_i \, \|f(z_i) - c_i\|_2^2.
\]
The repository implements this with sklearn regressors and `sample_weight`.
Each training example gets one scalar weight $\hat\alpha_i$, shared across all
cost coordinates.

\section*{4. Algorithm}

\begin{algorithm}[h]
\caption{Cross-fitted decision-weight regression}
\begin{algorithmic}[1]
\Require Data $\{(z_i,c_i)\}_{i=1}^n$, predictor class $\mathcal{F}$, weight-regression class $\mathcal{H}$, folds $I_1,\dots,I_M$, mixing weight $\mu \in [0,1]$.
\State Fit pilot predictor $\hat f$ by ordinary least squares.
\For{$i=1,\dots,n$}
    \State Compute pilot prediction $\hat c_i \gets \hat f(z_i)$.
    \State Compute regret $\Reg_i \gets c_i^\top(x^*(\hat c_i) - x^*(c_i))$.
    \State Form raw weight $w_i^{\mathrm{raw}} \gets (1-\mu) + \mu \phi(\Reg_i)$.
\EndFor
\For{$m=1,\dots,M$}
        \State Fit $q^{(-m)}$ on $\{(z_i,w_i^{\mathrm{raw}}): i \notin I_m\}$.
    \For{each $i \in I_m$}
        \State $\hat\alpha_i \gets \clip_{[1-\mu,\,1-\mu+\mu W_{\max}]}(q^{(-m)}(z_i))$.
    \EndFor
\EndFor
\State Fit final predictor $\tilde f$ using weighted least squares with sample weights $\hat\alpha_i$.
\State \Return $\tilde f$.
\end{algorithmic}
\end{algorithm}

\section*{5. Concrete sklearn implementation}

The code path is:

\begin{enumerate}
\item Fit the pilot cost model with \texttt{get\_weighted\_predictors(..., weights=None)}.
\item Compute training regrets with \texttt{get\_regret(...)}.
\item Convert regrets to raw weights with \texttt{get\_raw\_context\_weights(...)}.
\item Cross-fit the alpha weights with \texttt{estimate\_context\_weights\_cross\_fitted(...)}.
\item Refit the final cost model with \texttt{get\_weighted\_predictors(...)} using the estimated alpha vector as \texttt{sample\_weight}.
\end{enumerate}

The default weight regressor is ridge regression with internal tuning:
\[
  \texttt{RidgeCV}(\alpha \in 10^{-4}, 10^{-3}, \dots, 10^4).
\]
The experiment script also exposes:

\begin{itemize}
\item linear regression with no tuning
\item random forest regression with a small \texttt{GridSearchCV} sweep
\end{itemize}

\section*{6. Why cross-fitting is used}

If the same data were used both to fit the weight regressor and to evaluate its
predicted alpha weights on the same examples, the resulting weights could
overfit the realized regret noise. Cross-fitting avoids that direct reuse by
ensuring each $\hat\alpha_i$ is predicted out of fold.

\section*{7. Summary}

The implemented method is:

\begin{itemize}
\item decision-aware through regret-based raw weights
\item context-only through regression of those raw weights on $z$
\item regularized and tunable through sklearn weight regressors
\item minimal in code because the final update remains an ordinary weighted
least-squares refit
\end{itemize}

\end{document}
